\documentclass[11pt, a4paper]{article}

% bfw-style.tex — gemeinsame Style-Definitionen für alle Handouts/Aufgabenhefte
% Voraussetzung: Kompilation mit xelatex (wegen fontspec)

\usepackage[ngerman]{babel}
\usepackage{geometry}
\geometry{a4paper, top=2.2cm, bottom=2.2cm, left=2.3cm, right=2.3cm}

\usepackage{fontspec}
% Fallback-Kette: Source Sans Pro -> Calibri -> Segoe UI -> Latin Modern Sans
% (Latin Modern ist immer mit MiKTeX vorhanden)
\IfFontExistsTF{Source Sans Pro}{%
  \setmainfont{Source Sans Pro}[Ligatures=TeX]%
}{\IfFontExistsTF{Calibri}{%
  \setmainfont{Calibri}[Ligatures=TeX]%
}{\IfFontExistsTF{Segoe UI}{%
  \setmainfont{Segoe UI}[Ligatures=TeX]%
}{%
  \setmainfont{Latin Modern Sans}[Ligatures=TeX]%
}}}
\IfFontExistsTF{Source Code Pro}{%
  \setmonofont{Source Code Pro}[Scale=0.92]%
}{\IfFontExistsTF{Consolas}{%
  \setmonofont{Consolas}[Scale=0.92]%
}{%
  \setmonofont{Latin Modern Mono}[Scale=0.92]%
}}

\usepackage{xcolor}
% Farbpalette: hell, freundlich, modern
\definecolor{bfwPrimary}{HTML}{0EA5A4}    % Petrol/Türkis - Primär
\definecolor{bfwPrimaryDark}{HTML}{0F766E} % dunkler Petrol für Headings
\definecolor{bfwAccent}{HTML}{F59E0B}     % Amber - Akzent
\definecolor{bfwSuccess}{HTML}{10B981}    % Grün - Erfolg / Lösung
\definecolor{bfwWarn}{HTML}{F97316}       % Orange - Achtung
\definecolor{bfwInk}{HTML}{0F172A}        % fast Schwarz
\definecolor{bfwInkSoft}{HTML}{475569}    % gedämpftes Grau
\definecolor{bfwBgSoft}{HTML}{F8FAFC}     % off-white
\definecolor{bfwBgTeal}{HTML}{ECFEFF}     % helles Türkis
\definecolor{bfwBgAmber}{HTML}{FEF3C7}    % helles Gelb
\definecolor{bfwBgRose}{HTML}{FFE4E6}     % helles Rosé für Eselsbrücken

\usepackage[colorlinks=true, linkcolor=bfwPrimaryDark, urlcolor=bfwPrimaryDark]{hyperref}
\usepackage{enumitem}
\setlist[itemize]{leftmargin=*, itemsep=2pt, topsep=2pt}
\setlist[enumerate]{leftmargin=*, itemsep=2pt, topsep=2pt}

\usepackage[most]{tcolorbox}
\usepackage{booktabs}
\usepackage{tikz}
\usetikzlibrary{decorations.pathmorphing, positioning, shapes.geometric, arrows.meta}

% Merkkästchen — wichtige Definition / Kernpunkt
\newtcolorbox{merksatz}[1][]{%
  enhanced, breakable,
  colback=bfwBgTeal,
  colframe=bfwPrimary,
  boxrule=0.6pt,
  arc=4pt,
  left=10pt, right=10pt, top=8pt, bottom=8pt,
  fonttitle=\bfseries\color{bfwPrimaryDark},
  title={\faLightbulb\ Merksatz},
  #1
}

% Eselsbrücke — Mnemoniks
\newtcolorbox{eselsbruecke}[1][]{%
  enhanced, breakable,
  colback=bfwBgRose,
  colframe=bfwAccent,
  boxrule=0.6pt,
  arc=4pt,
  left=10pt, right=10pt, top=8pt, bottom=8pt,
  fonttitle=\bfseries\color{bfwWarn},
  title={Eselsbrücke},
  #1
}

% Beispiel-Box
\newtcolorbox{beispiel}[1][]{%
  enhanced, breakable,
  colback=bfwBgSoft,
  colframe=bfwInkSoft,
  boxrule=0.4pt,
  arc=3pt,
  left=10pt, right=10pt, top=8pt, bottom=8pt,
  fonttitle=\bfseries\color{bfwInk},
  title={Beispiel},
  #1
}

% Lösungs-Box (für Aufgabenhefte)
\newtcolorbox{loesung}[1][]{%
  enhanced, breakable,
  colback=bfwBgAmber!40,
  colframe=bfwSuccess,
  boxrule=0.6pt,
  arc=3pt,
  left=10pt, right=10pt, top=8pt, bottom=8pt,
  fonttitle=\bfseries\color{bfwSuccess!70!black},
  title={Lösung},
  #1
}

% Glossar-Eintrag
\newcommand{\glossareintrag}[2]{%
  \par\noindent\textbf{\color{bfwPrimaryDark}#1} \enspace #2\par\smallskip
}

% Section-Stil — etwas farbiger als Standard
\usepackage{titlesec}
\titleformat{\section}
  {\Large\bfseries\color{bfwPrimaryDark}}
  {\thesection}{0.6em}{}
\titleformat{\subsection}
  {\large\bfseries\color{bfwPrimary}}
  {\thesubsection}{0.5em}{}
\titleformat{\subsubsection}
  {\normalsize\bfseries\color{bfwInk}}
  {\thesubsubsection}{0.4em}{}

% TikZ-Stil "skizzenhaft" — etwas weicher als Rough.js, aber stabil
\tikzset{
  roughbox/.style = {
    draw=bfwPrimaryDark, thick, fill=bfwBgTeal,
    rounded corners=4pt, inner sep=6pt
  },
  roughaccent/.style = {
    draw=bfwAccent, thick, fill=bfwBgAmber,
    rounded corners=4pt, inner sep=6pt
  },
  roughline/.style = {
    draw=bfwInkSoft, thick, -{Stealth[length=2.5mm]}
  }
}

% Bullet-Symbole (bewusst kein FontAwesome — vermeidet Font-Installations-Probleme)
\newcommand{\faLightbulb}{$\bullet$}
\newcommand{\faBookmark}{$\bullet$}

% Header/Footer
\usepackage{fancyhdr}
\pagestyle{fancy}
\fancyhf{}
\renewcommand{\headrulewidth}{0pt}
\fancyfoot[L]{\small\color{bfwInkSoft}\thepage}
\fancyfoot[R]{\small\color{bfwInkSoft} BFW M-064 Beschaffung im Büromanagement}


\title{\vspace*{-1.5cm}\Huge\bfseries\color{bfwPrimaryDark}Aufgabenheft Tag~1\\[0.4em]
  \large\color{bfwInkSoft}\textnormal{Beschaffungsprozess --- Übungen im IHK-Klausurformat}}
\author{\textsf{Beschaffung im Büromanagement}\\
\textsf{\small\copyright\ Can Siebert\,\textperiodcentered\,2026}}
\date{Selbstlernmaterial \textperiodcentered{} 11.05.2026}

\begin{document}
\maketitle
\thispagestyle{empty}

\begin{merksatz}[title={\bfseries So nutzen Sie dieses Aufgabenheft}]
Bearbeiten Sie alle Aufgaben zunächst \emph{ohne} in die Lösungen zu schauen. Schreiben
Sie Ihre Antworten in vollständigen Sätzen --- die IHK-Prüfung verlangt formulierte
Begründungen, nicht bloße Stichworte. Beachten Sie die \emph{Operatoren} (nennen,
erläutern, beurteilen, begründen) --- sie geben den geforderten Antwort-Tiefegrad vor.
Erst danach öffnen Sie den Lösungsteil ab Seite~\pageref{loesungen} und vergleichen.
\end{merksatz}

\tableofcontents
\vspace{1em}
\hrule

\section{Aufgaben}

\subsection*{Aufgabe~1 --- Wertschöpfungskette und Beschaffung (10 Punkte)}

\noindent\textbf{\color{bfwAccent}YouTube-Vorbereitung:}\enspace \href{https://www.youtube.com/watch?v=Ta25vVbxNv0}{Wertschöpfungskette --- Was ist das?} (\url{https://www.youtube.com/watch?v=Ta25vVbxNv0})

\medskip
Sie sind seit drei Monaten in der Einkaufsabteilung der \emph{Schreibwaren-Vertrieb GmbH}
in Berlin tätig. Ihre Ausbilderin bittet Sie, einer neuen Auszubildenden den Begriff
\emph{Wertschöpfungskette} und die Rolle der Beschaffung darin zu erklären.

\begin{enumerate}[label=(\alph*)]
  \item \textbf{Erläutern} Sie in eigenen Worten den Begriff \emph{Wertschöpfungskette}. \hfill (3~P.)
  \item \textbf{Ordnen} Sie die Beschaffung innerhalb der Wertschöpfungskette eines Industrieunternehmens \textbf{ein}. \hfill (2~P.)
  \item \textbf{Nennen} Sie drei konkrete Auswirkungen, die eine fehlerhafte Beschaffung auf nachgelagerte Bereiche (Lager, Produktion, Vertrieb, Kundenservice) haben kann. \hfill (3~P.)
  \item \textbf{Beurteilen} Sie kurz, warum die Beschaffung in modernen Unternehmen zunehmend strategische Bedeutung gewonnen hat. \hfill (2~P.)
\end{enumerate}

\subsection*{Aufgabe~2 --- Die sechs Beschaffungsziele (12 Punkte)}

\noindent\textbf{\color{bfwAccent}YouTube-Vorbereitung:}\enspace \href{https://www.youtube.com/watch?v=DhvdMbDG0Xk}{Beschaffungsplanung --- einfach erklärt} (\url{https://www.youtube.com/watch?v=DhvdMbDG0Xk})

\medskip
Die ,,6~R'' der Beschaffung lauten: das richtige Material in der richtigen Menge in der
richtigen Qualität zum richtigen Zeitpunkt am richtigen Ort zum richtigen Preis.

\begin{enumerate}[label=(\alph*)]
  \item \textbf{Nennen} Sie alle sechs Beschaffungsziele in der korrekten Reihenfolge. \hfill (3~P.)
  \item \textbf{Geben} Sie zu jedem Ziel ein konkretes Praxisbeispiel aus dem Büroalltag \textbf{an}. \hfill (3~P.)
  \item \textbf{Erläutern} Sie zwei mögliche Zielkonflikte zwischen den 6~R und beschreiben Sie, wie Sie als Sachbearbeiter/-in damit umgehen würden. \hfill (4~P.)
  \item \textbf{Beurteilen} Sie die Aussage: \emph{,,Der billigste Lieferant ist immer die wirtschaftlichste Wahl.''} \hfill (2~P.)
\end{enumerate}

\subsection*{Aufgabe~3 --- Make or Buy: Reinigungsdienst (15 Punkte)}

\noindent\textbf{\color{bfwAccent}YouTube-Vorbereitung:}\enspace \href{https://www.youtube.com/watch?v=R_GnBQj68L4}{Make-or-Buy --- Berechnung erklärt} (\url{https://www.youtube.com/watch?v=R_GnBQj68L4})

\medskip
Die \emph{Müller GmbH} (mittelständischer Bürodienstleister, 80~Mitarbeitende) erwägt,
die tägliche Büroreinigung an einen externen Dienstleister zu vergeben. Aktuell
beschäftigt sie zwei eigene Reinigungskräfte. Die Geschäftsführung legt Ihnen folgende
Zahlen vor und bittet Sie um eine Entscheidungsempfehlung:

\begin{center}
\begin{tabular}{p{6cm} r r}
\toprule
\textbf{Position} & \textbf{Eigenleistung (€/Mon.)} & \textbf{Fremdbezug (€/Mon.)}\\
\midrule
Personalkosten Reinigung & 4\,200,00 & --- \\
Reinigungsmaterial & 350,00 & --- \\
Anteil Verwaltungskosten & 280,00 & --- \\
Externer Dienstleister (Pauschalvertrag) & --- & 4\,500,00 \\
\bottomrule
\end{tabular}
\end{center}

\begin{enumerate}[label=(\alph*)]
  \item \textbf{Berechnen} Sie die Gesamtkosten beider Varianten und ermitteln Sie die monatliche bzw.\ jährliche Kostendifferenz. \hfill (4~P.)
  \item \textbf{Nennen} Sie drei qualitative Faktoren, die für oder gegen den Fremdbezug sprechen. \hfill (3~P.)
  \item \textbf{Erläutern} Sie konkret das Risiko der \emph{Lieferantenabhängigkeit} bei einem Outsourcing-Vertrag. \hfill (3~P.)
  \item \textbf{Begründen} Sie, welche Empfehlung Sie der Geschäftsführung geben. Beziehen Sie quantitative und qualitative Argumente ein. \hfill (5~P.)
\end{enumerate}

\subsection*{Aufgabe~4 --- E-Procurement: Konzepte zuordnen (8 Punkte)}

\noindent\textbf{\color{bfwAccent}YouTube-Vorbereitung:}\enspace \href{https://www.youtube.com/watch?v=zIN1FewmZCU}{E-Procurement: Buy Side, Sell Side, Marktplatz} (\url{https://www.youtube.com/watch?v=zIN1FewmZCU})

\medskip
\begin{enumerate}[label=(\alph*)]
  \item \textbf{Nennen} und \textbf{erläutern} Sie kurz die drei Grundkonzepte des E-Procurement. \hfill (3~P.)
  \item \textbf{Ordnen} Sie folgende Beispiele dem passenden Konzept zu \textbf{und begründen} Sie Ihre Zuordnung jeweils in einem Satz: \hfill (3~P.)
    \begin{itemize}[leftmargin=*]
      \item Eine Mitarbeiterin bestellt Toner im Onlineshop von \emph{Office-Versand AG}.
      \item Der Konzern \emph{Bayer AG} betreibt ein eigenes Bestellportal, in dem zugelassene Lieferanten ihre Kataloge einstellen.
      \item Das Portal \emph{Mercateo} bündelt Sortimente vieler Lieferanten.
    \end{itemize}
  \item \textbf{Beurteilen} Sie kurz: Welcher Vorteil spricht für ein Unternehmen mit 50~Mitarbeitenden für die Nutzung eines \emph{Marktplatz}-Modells? \hfill (2~P.)
\end{enumerate}

\subsection*{Aufgabe~5 --- Ökologische Beschaffung und LkSG (12 Punkte)}

\noindent\textbf{\color{bfwAccent}YouTube-Vorbereitung:}\enspace \href{https://www.youtube.com/watch?v=S2pVrFUND-c}{LkSG kurz erklärt} (\url{https://www.youtube.com/watch?v=S2pVrFUND-c})

\medskip
Ein Lieferant Ihres Betriebs steht in der Kritik, weil er Vorprodukte aus einem Land
mit fragwürdigen Arbeitsbedingungen bezieht. Ihre Vorgesetzte fragt Sie, ob das für die
Müller GmbH (700~Mitarbeitende) rechtlich relevant ist und welche Bewertungs\-kriterien
in der Lieferantenauswahl ergänzt werden sollten.

\begin{enumerate}[label=(\alph*)]
  \item \textbf{Erläutern} Sie kurz das \emph{Lieferketten\-sorgfalts\-pflichten\-gesetz (LkSG)}: Was ist sein Ziel und für welche Unternehmen gilt es? \hfill (3~P.)
  \item \textbf{Beurteilen} Sie, ob die Müller GmbH (700~Mitarbeitende) direkt unter das LkSG fällt. Begründen Sie. \hfill (2~P.)
  \item \textbf{Nennen} Sie fünf ökologische bzw.\ soziale Kriterien, die in einer Lieferantenauswahl berücksichtigt werden können. \hfill (3~P.)
  \item \textbf{Begründen} Sie: Auch ohne direkte Gesetzespflicht --- warum sollte die Müller GmbH dennoch nachhaltige Beschaffungs\-kriterien anwenden? \hfill (4~P.)
\end{enumerate}

\subsection*{Aufgabe~6 --- Beschaffungsprinzip wählen (10 Punkte)}

\noindent\textbf{\color{bfwAccent}YouTube-Vorbereitung:}\enspace \href{https://www.youtube.com/watch?v=cDwQvphoaFU}{Beschaffungsarten einfach erklärt} (\url{https://www.youtube.com/watch?v=cDwQvphoaFU})

\medskip
\begin{enumerate}[label=(\alph*)]
  \item \textbf{Nennen} und \textbf{charakterisieren} Sie kurz die drei Beschaffungs\-prinzipien (Einzel-, Vorrats-, fertigungssynchrone Beschaffung). \hfill (3~P.)
  \item \textbf{Ordnen} Sie für jeden Fall das passende Beschaffungs\-prinzip zu \textbf{und begründen} Sie: \hfill (4~P.)
    \begin{enumerate}[label=\arabic*., leftmargin=*]
      \item Eine Druckerei benötigt täglich 200\,kg frische Papierrollen für die laufende Produktion.
      \item Ein Anwaltsbüro kauft jährlich 500 Aktenordner für das Vertragsarchiv ein.
      \item Eine Werft fertigt eine Sondermaschine zur Verladung von Containern; diese wird nur einmal benötigt.
      \item Ein Automobil\-zulieferer beliefert das Werk seines Kunden mit just-in-time gelieferten Karosserieteilen.
    \end{enumerate}
  \item \textbf{Beurteilen} Sie das größte \emph{Risiko} der Just-in-Time-Beschaffung anhand eines konkreten Beispiels. \hfill (3~P.)
\end{enumerate}

\vspace{1em}
\noindent\textbf{Punkte gesamt: 67}

\newpage

\section{Lösungen}\label{loesungen}

\subsection*{Lösung Aufgabe~1}
\begin{loesung}
\textbf{(a)} Die Wertschöpfungskette beschreibt die geordnete Abfolge betrieblicher
Tätigkeiten, durch die ein Unternehmen Inputs (Material, Energie, Wissen, Arbeitskraft)
schrittweise in verkaufsfertige Produkte oder Dienstleistungen umwandelt. Jeder Schritt
erhöht den Wert des Outputs --- vom eingekauften Rohmaterial bis zum betreuten Endkunden.

\textbf{(b)} Die Beschaffung steht ganz am \emph{Anfang} der Wertschöpfungskette. Sie
ist das Bindeglied zwischen dem Beschaffungsmarkt (Lieferanten) und der internen
Wertschöpfung (Lager, Produktion, Vertrieb).

\textbf{(c)} Drei beispielhafte Auswirkungen einer fehlerhaften Beschaffung:
\begin{itemize}[itemsep=0pt]
  \item \emph{Produktion}: Stillstand wegen fehlenden oder qualitativ unzureichenden Materials, was zu Termin\-verzögerungen und Mehrkosten führt.
  \item \emph{Vertrieb}: Lieferverzug an Kunden, dadurch Umsatzeinbußen und Vertragsstrafen.
  \item \emph{Kundenservice}: Reklamations\-aufkommen steigt, Beschwerden lösen Image\-verlust und Kundenabwanderung aus.
\end{itemize}

\textbf{(d)} Die strategische Bedeutung wächst, weil (1)~der Materialkostenanteil am
Endprodukt bei vielen Unternehmen über 50\,\% liegt --- jede Einsparung schlägt direkt
auf den Gewinn durch, (2)~globalisierte Lieferketten Risiken verstärken (Lieferanten\-
ausfall, geopolitische Krisen) und (3)~Nachhaltigkeit\-spflichten (LkSG, CSRD)
strategisches Lieferantenmanagement zur gesetzlichen Anforderung machen.
\end{loesung}

\subsection*{Lösung Aufgabe~2}
\begin{loesung}
\textbf{(a)} Die sechs Beschaffungsziele in korrekter Reihenfolge:
\textbf{Material, Menge, Qualität, Ort, Zeit, Preis} (Eselsbrücke: M-M-Q-O-Z-P).

\textbf{(b)} Praxisbeispiele aus dem Büroalltag:
\begin{itemize}[itemsep=0pt]
  \item \emph{Material}: ,,Aktenordner DIN~A4 breit, blau, mit Hebel\-mechanik'' --- nicht nur ,,Aktenordner''.
  \item \emph{Menge}: 50 Aktenordner für das Quartal --- nicht 1\,000 wegen Mengenrabatt überbestellen.
  \item \emph{Qualität}: 80~g/m² Druckerpapier für Bürodruck --- nicht Premium-Fotopapier.
  \item \emph{Ort}: Lieferung in das Archivlager (3.\,OG) --- nicht in den Wareneingang.
  \item \emph{Zeit}: Lieferung am vereinbarten Werktag --- nicht 14~Tage zu früh.
  \item \emph{Preis}: Bezugspreis (inkl.\ Bezugskosten) als Vergleichsmaßstab --- nicht der Listenpreis.
\end{itemize}

\textbf{(c)} Beispielhafte Zielkonflikte:
\begin{itemize}[itemsep=0pt]
  \item \emph{Menge gegen Lagerkosten}: Mengenrabatt senkt den Stückpreis, erhöht aber Lagerkosten und Kapital\-bindung. \emph{Umgang}: Bestellmengen über die optimale Bestellmenge (Andler-Formel) wirtschaftlich ermitteln, nicht emotional.
  \item \emph{Preis gegen Qualität}: Billiger Lieferant verursacht oft höhere Folgekosten durch Reklamationen und Liefer\-verzug. \emph{Umgang}: Lieferanten\-bewertung mit Nutzwertanalyse, nicht nur Preisvergleich.
\end{itemize}

\textbf{(d)} Die Aussage ist falsch. Der billigste Lieferant ist nicht automatisch der wirtschaftlichste, weil Folgekosten (Reklamationen, Liefer\-verzug, Image\-schäden) nicht im Listenpreis enthalten sind. Wirtschaftlich ist derjenige Lieferant, der \emph{insgesamt} (inkl.\ aller Bezugs- und Folgekosten) die beste Gesamt\-bilanz liefert.
\end{loesung}

\subsection*{Lösung Aufgabe~3}
\begin{loesung}
\textbf{(a)} Quantitativer Vergleich:

\begin{center}
\begin{tabular}{p{5cm} r r}
\toprule
 & Eigenleistung & Fremdbezug\\
\midrule
Personalkosten & 4\,200,00\,€ & --- \\
Reinigungsmaterial & 350,00\,€ & --- \\
Verwaltungskosten & 280,00\,€ & --- \\
Externer Dienstleister & --- & 4\,500,00\,€ \\
\midrule
\textbf{Gesamt / Monat} & \textbf{4\,830,00\,€} & \textbf{4\,500,00\,€}\\
\midrule
\textbf{Differenz / Monat} & \multicolumn{2}{r}{\textbf{330,00\,€ zugunsten Fremdbezug}}\\
\textbf{Differenz / Jahr} & \multicolumn{2}{r}{\textbf{3\,960,00\,€ ($\approx$\,6{,}8\,\%)}}\\
\bottomrule
\end{tabular}
\end{center}

\textbf{(b)} Drei qualitative Faktoren:
\begin{itemize}[itemsep=0pt]
  \item \emph{Qualitätskontrolle}: Bei Eigenleistung direkt kontrollierbar, bei Fremdbezug nur indirekt über Vertragsmanagement und Service Level Agreements (SLA).
  \item \emph{Flexibilität}: Eigenes Personal reagiert spontan auf Sonderwünsche (z.\,B.\ Reinigung nach Kunden\-besuch). Externe Dienstleister sind an Vertragsstandards gebunden.
  \item \emph{Sozialer Faktor}: Eigene Mitarbeitende kennen das Haus, die Abläufe und die Kollegen. Wechselnde Fremdkräfte sind anonymer.
\end{itemize}

\textbf{(c)} \emph{Lieferantenabhängigkeit} entsteht, sobald keine eigene Kapazität mehr vorhanden ist. Konkretes Risiko: Erhöht der Dienstleister nach 6~Monaten den Preis um 20\,\%, hat die Müller GmbH kaum Verhandlungsmacht --- denn die eigenen Reinigungskräfte sind entlassen, der Wieder\-aufbau eigener Kapazität würde Wochen dauern. Der Dienstleister kann diese Wechselkosten ausnutzen (\emph{Konditionsmacht}).

\textbf{(d)} Empfehlung: Der quantitative Vorteil von 6{,}8\,\% ist gering. Bei dieser Größen\-ordnung überwiegen häufig die qualitativen Risiken. Empfehlung an die Geschäftsführung:
\begin{itemize}[itemsep=0pt]
  \item Wenn Outsourcing, dann mit \emph{klar definiertem SLA} (Reinigungs\-zeiten, Qualitäts\-kontrollen, Pönalen bei Mängeln).
  \item Kurze Vertrags\-laufzeit von max.\ 12 Monaten mit Sonder\-kündigungs\-recht.
  \item Übergangs\-management: bestehende Reinigungs\-kräfte angemessen abfindigen oder versetzen.
  \item \emph{Rückfall-Option} sicherstellen: dokumentierter Reinigungs\-prozess, der bei Bedarf wieder eigenes Personal anlernen kann.
\end{itemize}
Bei Erfüllung dieser Bedingungen ist der Fremdbezug vertretbar; ohne diese Sicherungen wäre die Eigenleistung trotz höherer Kosten die robustere Wahl.
\end{loesung}

\subsection*{Lösung Aufgabe~4}
\begin{loesung}
\textbf{(a)} Die drei E-Procurement-Konzepte:
\begin{itemize}[itemsep=0pt]
  \item \emph{Sell-Side-Lösung}: Der Lieferant betreibt das System (z.\,B.\ Onlineshop). Käufer kaufen dort ein.
  \item \emph{Buy-Side-Lösung}: Der Käufer betreibt das System; zugelassene Lieferanten stellen ihre Kataloge ein. Typisch bei Konzernen mit zentralem Einkauf.
  \item \emph{Marktplatz}: Ein neutraler Dritter bündelt Sortimente vieler Lieferanten und stellt die Plattform bereit.
\end{itemize}

\textbf{(b)} Zuordnung mit Begründung:
\begin{itemize}[itemsep=0pt]
  \item \emph{Office-Versand-Onlineshop}: \textbf{Sell-Side} --- der Lieferant betreibt das System.
  \item \emph{Bayer-Bestellportal}: \textbf{Buy-Side} --- der Käufer (Bayer) betreibt das System; Lieferanten haben Lese-/Einstell\-rechte.
  \item \emph{Mercateo}: \textbf{Marktplatz} --- ein neutraler Dritter bündelt viele Anbieter.
\end{itemize}

\textbf{(c)} Vorteil eines Marktplatzes für ein Unternehmen mit 50~Mitarbeitenden: Eine eigene Buy-Side-Lösung wäre wirtschaftlich nicht darstellbar (Investitions- und Pflege\-kosten zu hoch). Der Marktplatz bietet eine breite Lieferanten\-auswahl und einen einheitlichen Bestell\-workflow ohne IT-Investitionen --- ideale Lösung für KMU.
\end{loesung}

\subsection*{Lösung Aufgabe~5}
\begin{loesung}
\textbf{(a)} Das \emph{Lieferketten\-sorgfalts\-pflichten\-gesetz (LkSG)} ist ein deutsches
Bundes\-gesetz, das seit dem 1.~Januar 2023 in Kraft ist. Es verpflichtet Unternehmen
einer bestimmten Größe, Sorgfalts\-pflichten in ihrer Lieferkette wahrzunehmen ---
insbesondere im Hinblick auf den Schutz von Menschenrechten (Verbot von Kinderarbeit,
faire Arbeits\-bedingungen, Schutz vor Diskriminierung) und der Umwelt entlang der
gesamten Wert\-schöpfungs\-kette. Der Geltungsbereich umfasste 2023 Unternehmen ab
3\,000 Mitarbeitenden und wurde seit 1.~Januar 2024 auf Unternehmen ab 1\,000~Mit\-arbeitenden
ausgeweitet. Die Aufsicht führt das Bundesamt für Wirtschaft und Ausfuhrkontrolle (BAFA);
Verstöße können mit Bußgeldern bis 8\,Mio.\,€ oder 2\,\% des Jahres\-umsatzes geahndet werden.

\textbf{(b)} Die Müller GmbH mit 700 Mitarbeitenden fällt \emph{nicht direkt} unter das
LkSG, da der Schwellenwert von 1\,000 Mitarbeitenden unterschritten wird. \emph{Indirekt}
ist sie aber sehr wohl betroffen: Wenn die Müller GmbH zur Lieferkette eines LkSG-pflichtigen
Großunternehmens gehört, wird dieses von ihr Sorgfalts\-nachweise verlangen. Die
geplante EU-Lieferketten\-richtlinie (CSDDD) wird zudem die Schwellen weiter senken.

\textbf{(c)} Fünf Bewertungs\-kriterien für nachhaltige Lieferanten:
\begin{itemize}[itemsep=0pt]
  \item Transport\-wege (möglichst kurz, regional)
  \item Verpackungs\-minimierung und Recyclingfähigkeit
  \item Energie\-bilanz und CO\textsubscript{2}-Fußabdruck
  \item Sozialstandards entlang der gesamten Lieferkette (Faire Löhne, sichere Arbeit)
  \item Zertifikate (z.\,B.\ ISO 14001 für Umweltmanagement, SA 8000 für Sozial\-standards)
\end{itemize}

\textbf{(d)} Auch ohne direkte Gesetzespflicht sollte die Müller GmbH nachhaltig
beschaffen, weil:
\begin{itemize}[itemsep=0pt]
  \item \emph{Indirekte Pflicht}: Großkunden verlangen Sorgfalts\-nachweise --- sonst droht der Auftrags\-verlust.
  \item \emph{Reputation}: Skandale entlang der Lieferkette schaden dem Image und damit dem Umsatz.
  \item \emph{Risiko\-vermeidung}: Lieferanten mit fairen Arbeits\-bedingungen und stabilen Strukturen sind verlässlicher.
  \item \emph{Künftige Gesetzes\-pflicht}: CSDDD und CSRD werden den Geltungsbereich erweitern --- frühzeitige Vorbereitung spart spätere Hektik.
\end{itemize}
\end{loesung}

\subsection*{Lösung Aufgabe~6}
\begin{loesung}
\textbf{(a)} Die drei Beschaffungs\-prinzipien:
\begin{itemize}[itemsep=0pt]
  \item \emph{Einzelbeschaffung im Bedarfsfall}: Pro konkretem Auftrag wird gezielt eingekauft. Kein Lager, hoher administrativer Aufwand. Geeignet für einmalige oder seltene Bedarfe.
  \item \emph{Vorratsbeschaffung}: Größere Mengen werden auf Lager gelegt. Hohe Versorgungssicherheit, hohe Kapital\-bindung. Geeignet für Standard\-verbrauchsmaterial mit gleichmäßigem Bedarf.
  \item \emph{Fertigungssynchrone Beschaffung (Just-in-Time)}: Lieferung erfolgt exakt zum Bedarfszeit\-punkt. Lager praktisch null, hohe Liefer\-anten\-abhängigkeit. Geeignet für gleichmäßigen, planbaren Bedarf mit zuverlässigen Lieferanten.
\end{itemize}

\textbf{(b)} Zuordnung mit Begründung:
\begin{enumerate}[label=\arabic*., itemsep=0pt]
  \item Druckerei mit täglichem Papierbedarf: \textbf{Just-in-Time}. Der gleichmäßige tägliche Bedarf rechtfertigt fertigungssynchrone Lieferung mit minimalem Lager.
  \item Anwaltsbüro mit jährlich 500 Aktenordnern: \textbf{Vorratsbeschaffung}. Standardartikel mit gleichmäßigem Verbrauch --- größere Bestellung mit Mengenrabatt, kleines Lager.
  \item Werft mit Sondermaschine: \textbf{Einzelbeschaffung im Bedarfsfall}. Einmaliger Bedarf, keine Wiederbeschaffung --- gezielter Einzelkauf.
  \item Automobil\-zulieferer mit Karosserieteilen: \textbf{Just-in-Time}. Klassische JIT-Branche; enge IT-Integration mit dem Kunden, gemeinsame Bestandsdaten.
\end{enumerate}

\textbf{(c)} Größtes Risiko der Just-in-Time-Beschaffung: \emph{Lieferketten\-störung}.
Konkretes Beispiel --- die Halbleiter\-knappheit 2021/2022: Viele Automobil\-werke mussten
ihre Produktion tagelang oder wochenlang stoppen, weil JIT-gelieferte Mikro\-chips
ausblieben. Da kein Sicherheits\-bestand vorhanden war, wurde aus dem Liefer\-engpass
sofort ein Produktions\-stopp mit Millionen\-verlusten. JIT senkt also Lagerkosten
massiv, erhöht aber die System\-anfälligkeit für externe Schocks dramatisch.
\end{loesung}

\end{document}
